% !TEX root = ../FundationsDataScience.tex

\chapter{Fourier and Convolution}
\label{sec-fourier}

\chapteroverview{Convolution combines shifted copies of a signal, but evaluating it directly can be costly and obscure the structure of the operation. Fourier coordinates turn convolution into multiplication, providing both efficient algorithms and a way to solve differential equations. We develop this connection on continuous and finite domains, study sampling and the fast Fourier transform, and extend the viewpoint to groups, surfaces, and graphs.}

The main reference for this chapter is \cite{mallat2008wavelet}.

\newcommand{\piN}{\frac{2\imath\pi}{N}}

%%%%%%%%%%%%%%%%%%%%%%%%%%%%%%%%%%%%%%%%%%%%%%%%
%%%%%%%%%%%%%%%%%%%%%%%%%%%%%%%%%%%%%%%%%%%%%%%%
%%%%%%%%%%%%%%%%%%%%%%%%%%%%%%%%%%%%%%%%%%%%%%%%
\section{Hilbert Spaces and Fourier Transforms}

%%%%%%%%%%%%%%%%%%%%%%%%%%%%%%%%%%%%%%%%%%%%%%%%
\subsection{Orthonormal Bases}

Many methods in data science begin by expanding the input signal in a basis. An orthonormal basis gives a simple reconstruction formula and preserves energy, which simplifies the analysis. We first describe these properties in a Hilbert space, such as $\Hh=L^2([0,1]^d)$ for signals on a continuous domain or $\Hh=\RR^N$ for discrete signals.

A complex Hilbert space $(\Hh,\dotp{\cdot}{\cdot})$ is complete for the norm induced by its Hermitian inner product. We take this inner product to be linear in the first argument and conjugate-linear in the second. A separable Hilbert space admits an orthonormal basis $(\phi_k)_k$ (finite in finite dimension), so every $f\in\Hh$ has the expansion
\eq{
	f = \sum_{k} \dotp{f}{\phi_k} \phi_k 
}
where the norm is determined by $\norm{f}^2 \eqdef \dotp{f}{f}$. Convergence means that $\norm{ f - \sum_{k=0}^N \dotp{f}{\phi_k} \phi_k } \rightarrow 0$ as $N \rightarrow +\infty$. Parseval\textquotesingle s identity expresses conservation of energy:
\eq{
	\norm{f}^2 = \sum_k |\dotp{f}{\phi_k}|^2.
}

The Gram--Schmidt procedure constructs such a basis from a linearly independent family $(\bar\phi_k)_k$ with dense span. Set $\phi_0=\bar\phi_0/\norm{\bar\phi_0}$ and $\phi_k=\tilde\phi_k/\norm{\tilde\phi_k}$, where $\tilde\phi_k$ is the orthogonal residual $\tilde\phi_k=\bar\phi_k - \sum_{i<k} a_i \phi_i$. Orthogonality requires $a_i = \dotp{\bar \phi_k}{\phi_i}$.

On $L^2([-1,1])$ with the usual inner product, orthogonalizing the monomials gives the Legendre polynomials. Their conventional normalization differs from unit $L^2$ norm:
\eq{
	\phi_0(x)=1, \quad \phi_1(x)=x, \quad \phi_2(x) = \frac{1}{2}(3x^2-1), \quad \text{etc.}
}
On $L^2(\RR,e^{-x^2/2}\,\d x)$, this construction gives the probabilists' Hermite polynomials $P_0(x)=1$, $P_1(x)=x$, $P_2(x)=x^2-1$, up to normalization. Multiplying normalized polynomials by $e^{-x^2/4}$ instead gives orthonormal Hermite functions for Lebesgue measure. A degree-$k$ orthogonal polynomial has exactly $k$ distinct real zeros in the interior of the support of the measure.

The Shannon interpolation theorem gives another example: $( \sinc(x-k) )_k$ is an orthonormal basis of the subspace $\enscond{f\in L^2(\RR)}{\supp(\hat f) \subset [-\pi,\pi]}$. Using the continuous representative of $f$, the reconstruction formula $f=\sum_k f(k) \sinc(x-k)$ identifies the coefficients as $\dotp{f}{\sinc(x-k)} = f(k)$. The series converges in $L^2$ and pointwise; the previous chapter proves the pointwise formula under additional decay assumptions.


Orthogonal bases can also arise as complete families of eigenvectors of self-adjoint operators. Below, we explain how translation-invariant operators (convolutions) characterize Fourier bases.


%%%%%%%%%%%%%%%%%%%%%%%%%%%%%%%%%%%%%%%%%%%%%%%%
\subsection{Fourier Basis on \texorpdfstring{$\RR/2\pi\ZZ$}{the Circle}}

Fourier analysis takes different forms on different domains. On the circle, it gives a countable orthonormal basis. On the real line, frequencies form a continuum, and the Fourier transform is not an expansion in a countable basis of exponential functions.

On $L^2(\TT)$ where $\TT=\RR/2\pi\ZZ$, equipped with $\dotp{f}{g} \eqdef \frac{1}{2\pi}\int_\TT f(x) \bar g(x) \d x$, one can use the Fourier basis 
\eql{\label{eq-fourier-basis-1d}
	\phi_k(x) \eqdef e^{\imath k x}
	\qforq
	k \in \ZZ. 
}
One thus has
\eql{\label{eq-defn-fourier-coeffs}
	f = \sum_{k\in\ZZ} \hat f_k e^{\imath k \cdot}
	\qwhereq
	\hat f_k \eqdef \frac{1}{2\pi}\int_0^{2\pi} f(x) e^{-\imath k x} \d x, 
}
with convergence in $L^2(\TT)$. Pointwise convergence is delicate; see Section~\ref{subsec-sampling}.

Figure \ref{fig-fourier-wav}, left, shows examples of the real parts of Fourier atoms.

\myfigure{
	\image{orthobases}{.4}{atoms-fourier}
	\image{orthobases}{.5}{atoms-wavelets}
}{%
	Examples of 1-D Fourier atoms (real parts, left) and wavelets (right). %
}{fig-fourier-wav}


We recall that for $f \in L^1(\RR)$, its Fourier transform is defined as
\eq{
	\foralls \om \in \RR, \quad
	\hat f(\om) \eqdef \int_\RR f(x) e^{-\imath x \om} \d x.
}
The transform extends to $L^2(\RR)$ by density.

The following diagram relates the Fourier transform on $\RR$ to Fourier coefficients on $\TT$:
\eq{
	\begin{array}{rcccl}
						& f(x)  &  \overset{\Ff}{\longrightarrow} &  \hat f(\om) &\\
		\text{sampling}& \downarrow & & \downarrow &\text{periodization} \\
						& (f(n))_n  &  \overset{\text{Fourier series}}{\longrightarrow} &  \sum_n f(n) e^{-\imath \om n} &\\
	\end{array}.
}
The two routes through the diagram give the identity
\eql{\label{eq-poisson-bis}
	\sum_n f(n) e^{-\imath \om n} = \sum_n \hat f(\om-2\pi n)
}
which is Poisson's summation formula (Proposition~\ref{prop-poisson}).



\myfigure{
\image{fourier}{.8}{fourier-transforms}
}{%
	Four settings for Fourier analysis, linked by sampling and periodization.%
}{fig-fourier-transforms}



%%%%%%%%%%%%%%%%%%%%%%%%%%%%%%%%%%%%%%%%%%%%%%%%
%%%%%%%%%%%%%%%%%%%%%%%%%%%%%%%%%%%%%%%%%%%%%%%%
%%%%%%%%%%%%%%%%%%%%%%%%%%%%%%%%%%%%%%%%%%%%%%%%
\section{Convolution on \texorpdfstring{$\RR$ and $\TT$}{the Line and the Circle}}

%%%%%%%%%%%%%%%%%%%%%%%%%%%%%%%%%%%%%%%%%%%%%%%%
\subsection{Convolution}

\wrapf{2-fourier/convolution}{Convolution on $\RR$.}
On $\XX=\RR$ or $\TT$, one defines 
\eql{\label{eq-convol-1d}
	f \star g(x) = \int_\XX f(t) g(x-t) \d t. 
}
Here and throughout this section, integrals use Lebesgue measure on $\RR$ and normalized Haar measure $\d t/(2\pi)$ on $\TT$. With this convention, Young's inequality states that, for $1\leq p,q,r\leq\infty$ and $(f,g)\in L^p(\XX)\times L^q(\XX)$,
\eql{\label{eq-young-convol}
	\frac{1}{p}+\frac{1}{q}=1+\frac{1}{r} \qarrq
	f \star g \in L^r(\XX) \qandq 
	\norm{f\star g}_{L^r(\XX)} \leq \norm{f}_{L^p(\XX)}\norm{g}_{L^q(\XX)}.
}
In particular, convolution by $f\in L^1(\XX)$ defines a bounded linear map from $L^p(\XX)$ to itself. When $r=\infty$ and $1<p,q<\infty$, the convolution $f\star g$ has a bounded continuous representative, illustrating its regularizing effect.
%
On the circle, $p<q \qarrq L^q(\TT) \subset L^p(\TT)$, so $L^\infty(\XX)$ is contained in every space in this scale.


\myfigure{
\image{fourier}{.6}{filtering-principle}
}{%
	Signal filtering with a box filter (running average). %	
}{fig-filter-box}

Convolution can regularize functions. For instance, if $f \in L^1(\XX)$ and $g\in C^1(\XX)$ is bounded with bounded derivative, then $f \star g$ is differentiable and $(f \star g)'=f \star g'$. To construct an approximate identity, choose a smooth compactly supported $\rho\geq0$ with $\int_\RR\rho=1$ and set $\rho_\epsilon(x)=\epsilon^{-1}\rho(x/\epsilon)$. Then $f\star\rho_\epsilon\to f$ in $L^p(\RR)$ for $1\leq p<\infty$, and uniformly for bounded uniformly continuous $f$. On the circle, use $2\pi\sum_{k\in\ZZ}\rho_\epsilon(\cdot-2\pi k)$, whose integral against normalized Haar measure is one.
%
This smoothing effect is also useful for denoising signals and images.


\myfigure{
\image{fourier}{.6}{filtering-1d-1}\\
\image{fourier}{.6}{filtering-1d-2}\\
\image{fourier}{.6}{filtering-1d-3}
}{%
	Filtering an irregular signal with Gaussian filters of increasing width $\si$. %
}{fig-filtering-1d}

%%%
\begin{prop}
If $f,g\in L^1(\XX)$, then
\eql{\label{eq-convol-fourier}
	\Ff(f \star g) = \hat f \odot \hat g.
}
For $\XX=\RR$, $\odot$ denotes pointwise multiplication of the continuous Fourier transforms.
For $\XX=\TT$, $\odot$ acts on bounded Fourier-coefficient sequences (pointwise multiplication of sequences).
\end{prop}


%%% FIG %%%
\wrapf{2-fourier/convolution-fourier}{How the Fourier transform converts convolution into multiplication.}
This means that $\Ff$ is an algebra homomorphism. For instance, if $\XX=\RR$, its range lies in the algebra of continuous functions that vanish at $\pm\infty$.

As shown in Figure~\ref{fig-splines}, successive convolutions of a box function produce cardinal splines: piecewise polynomials of increasing smoothness.

\begin{figure}
\centering
\includegraphics[width=.3\linewidth]{2-fourier/splines}
\caption{\label{fig-splines}
Cardinal splines are defined by successive convolutions.
}
\end{figure}


Convolution also describes the sum of independent random variables. Writing $f_X$ for the probability density of a random vector $X$ with respect to Lebesgue measure, independence of $X$ and $Y$ gives $f_{X+Y}=f_X \star f_Y$. This identity is fundamental in probability and statistics.


\texttt{Associated code: fourier/test\_denoising.m and fourier/test\_fft2.m}

%%%%%%%%%%%%%%%%%%%%%%%%%%%%%%%%%%%%%%%%%%%%%%%%
\subsection{Translation Invariant Operators}

\wrapf{2-fourier/translation-inv}{Commutative diagram for translation invariance.}
Translation-invariant operators commute with translations. In signal and image processing, they apply the same rule at every location.

The following propositions characterize translation-invariant operators as convolutions against kernels that may be distributions. The term ``translation equivariant'' is more precise: translating the input translates the output by the same amount.
%
The kernel\textquotesingle s regularity depends on the topologies of the input and output spaces, so the proof differs between the two settings below.
%
We begin with operators whose outputs are continuous functions.

\begin{prop}
	We define $T_\tau f = f(\cdot-\tau)$.
	%
	A bounded linear operator $H : (L^2(\XX),\norm{\cdot}_2) \rightarrow (C_b(\XX),\norm{\cdot}_\infty)$ commutes with every translation, meaning that for all $\tau$, $H \circ T_\tau = T_\tau \circ H$, if and only if
	\eq{
		\foralls f \in L^2(\XX), \quad H(f) = f \star g
	}
	with $g \in L^2(\XX)$.
\end{prop}

Formally, the identity $f=f\star\delta$ suggests that $Hf=f\star H\delta$ for a translation-invariant operator. Since $H\delta$ need not be defined under the present hypotheses, the following argument uses bounded evaluation functionals instead.

\begin{proof}
	By~\eqref{eq-young-convol} and the continuity result for $r=\infty$, the convolution operator $H : f \mapsto f \star g$ with $g \in L^2(\XX)$ is bounded from $L^2(\XX)$ to $C_b(\XX)$. Moreover,
	\eq{
		(H \circ T_\tau)(f) = \int_\XX f( (\cdot-\tau)-y) g(y) \d y = (f \star g)(\cdot-\tau) = T_\tau (Hf),
	}
	so that such an $H$ is translation-invariant. 
	
	
	Conversely, we define $\ell : f \mapsto H(f)(0) \in \CC$, which is well defined since $H(f) \in \Cc^0(\XX)$.
	%
	Continuity of $H$ gives a constant $C$ such that $\norm{Hf}_\infty \leq C \norm{f}_2$, hence $|\ell(f)| \leq C \norm{f}_2$. Thus $\ell$ is a continuous linear functional on the Hilbert space $L^2(\XX)$. By the Fr\'echet--Riesz theorem, there exists $h \in L^2(\XX)$ such that $\ell(f)=\dotp{f}{h}$.
	%
	Hence, $\foralls x \in \XX$, 
	\eq{
		H(f)(x) = T_{-x} (Hf)(0) = H ( T_{-x} f)(0) = \ell( T_{-x} f ) = \dotp{T_{-x} f}{h} =
		\int_\XX f(y+x)\overline{h(y)}\,\d y = (f\star g)(x),
	}
	where $g(x)\eqdef\overline{h(-x)}\in L^2(\XX)$.
\end{proof}

On $\TT$, an operator need not produce continuous outputs. Its convolution kernel may then be a distribution, and Fourier coefficients provide a convenient way to define the operator.

\begin{prop}\label{prop-ti-convol-l2}
	A bounded linear operator $H : L^2(\TT) \rightarrow L^2(\TT)$ commutes with every translation, meaning that for all $\tau$, $H \circ T_\tau = T_\tau \circ H$, if and only if
	\eq{
		\foralls f \in L^2(\TT), \quad
		\Ff( H(f) ) = \hat f \odot c 
	}
	where $c \in \ell^\infty(\ZZ)$ (a bounded sequence).
\end{prop}
\begin{proof}
	We denote $\phi_n \eqdef e^{\imath n \cdot}$. 
	One has
	\eq{
		H(\phi_n) = H(T_\tau T_{-\tau }\phi_n) = H(T_\tau e^{\imath n \tau } \phi_n)  = e^{\imath n \tau} H(T_\tau(\phi_n)) = e^{\imath n \tau} T_\tau(H(\phi_n)).
	}
	Thus, for all $n$, 
	\eq{
		\dotp{H(\phi_n)}{\phi_m} 
		= e^{\imath n \tau} \dotp{T_\tau \circ H(\phi_n)}{\phi_m} 
		= e^{\imath n \tau} \dotp{H(\phi_n)}{ T_{-\tau}(\phi_m) }
		=  e^{\imath (n-m) \tau} \dotp{H(\phi_n)}{ \phi_m }
	}
	So for $n \neq m$, $\dotp{H(\phi_n)}{\phi_m} = 0$, and we define $c_n \eqdef \dotp{H(\phi_n)}{\phi_n}$.
	We thus have $H(\phi_n) = c_n \phi_n$.
	Since $H$ is continuous, $\norm{Hf}_{L^2(\TT)} \leq C \norm{f}_{L^2(\TT)}$ for some constant $C$, and thus
	by Cauchy--Schwarz
	\eq{
		|c_n| = |\dotp{H(\phi_n)}{\phi_n}| \leq \norm{H(\phi_n)} \norm{\phi_n} \leq C 
	}	
	because $\norm{\phi_n}=1$, so that $c \in \ell^\infty(\ZZ)$.
	%
	By continuity of $H$, recalling that by definition $\hat f_n \eqdef \dotp{f}{\phi_n}$, 
	\eq{
		H(f) = \lim_N H( \sum_{n=-N}^N \hat f_n \phi_n ) = \lim_N \sum_{n=-N}^N \hat f_n H(\phi_n)
		= \lim_N \sum_{n=-N}^N c_n \hat f_n \phi_n = \sum_{n \in \ZZ} c_n \hat f_n \phi_n
	} 
	where we used $H(\phi_n) = c_n \phi_n$. This proves the desired representation.
\end{proof}

Conversely, by Parseval\textquotesingle s identity, a bounded sequence $c$ defines a bounded Fourier multiplier, which commutes with translations. Translation-invariant operators are therefore diagonal in the Fourier basis.

%%%%%%%%%%%%%%%%%%%%%%%%%%%%%%%%%%%%%%%%%%%%%%%%
\subsection{Poisson\textquotesingle s Formula and Distributions}

Informally, the Fourier series
\eq{
	\sum_n f(n) e^{-\imath \om n}
} 
can be viewed as the Fourier transform $\Ff( \Pi_1 \odot f )$ of the discrete distribution
\eq{
	\Pi_1 \odot f = \sum_n f(n) \de_{n}
	\qwhereq
	\Pi_s = \sum_n \de_{sn}
}
for $s>0$,
where $\de_a$ is the Dirac mass at location $a \in \RR$, i.e. the distribution such that $\int f\,\d\delta_a=f(a)$ for any continuous $f$. A measure can be multiplied by a continuous function; a general distribution can be multiplied by a smooth function. For a tempered distribution $\mu$, its Fourier transform is defined by duality on Schwartz test functions:
\eq{
	\foralls g \in \Ss(\RR), \quad \int_\RR g(x) \d \Ff(\mu) = \int_\RR \Ff(g) \d \mu,
	\qwhereq
	\Ff(g) \eqdef \int_\RR g(x)e^{-\imath x\cdot}\,\d x,
}
Here $\Ss(\RR)$ is the Schwartz space; the integral notation denotes the distributional pairing.

\wrapf{2-fourier/poisson-distrib}{Sine waves summed in the Poisson formula.}
The Poisson formula~\eqref{eq-poisson-bis} can thus be interpreted as
\eq{
	\Ff(\Pi_1 \odot f) = \sum_n \hat f(\cdot-2\pi n) = \int_\RR \hat f(\cdot-\om) \d\Pi_{2\pi}(\om)  = \hat f \star \Pi_{2\pi}
}
Since $\Ff^{-1} = \frac{1}{2\pi} \Ss \circ \Ff$ where $\Ss(f)=f(-\cdot)$, applying this operator to both sides gives
\eq{
	\Pi_1 \odot f = \frac{1}{2\pi} \Ss \circ \Ff(\hat f \star \Pi_{2\pi}) =  \Ss(  \frac{1}{2\pi} \Ff(\hat f) \odot \hat\Pi_{2\pi})
	= \Ss(  \frac{1}{2\pi} \Ff(\hat f) ) \odot \Ss(\hat\Pi_{2\pi}) = \hat \Pi_{2\pi} \odot f.
}
This can be interpreted as the relation
\eq{
	\hat\Pi_{2\pi} = \Pi_{1}
	\qarrq
	\hat\Pi_1 = 2\pi \Pi_{2\pi}.
}

For intuition, consider a finite Fourier series
\eq{
	\sum_{n=-N}^N e^{-\imath n \om} = \frac{\sin((N+1/2)\om)}{\sin(\om/2)}
}
which is the Dirichlet kernel, with removable singularities at multiples of $2\pi$. Its value there is $2N+1$; in the sense of distributions it converges to $2\pi\Pi_{2\pi}$.

%%%%%%%%%%%%%%%%%%%%%%%%%%%%%%%%%%%%%%%%%%%%%%%%
%%%%%%%%%%%%%%%%%%%%%%%%%%%%%%%%%%%%%%%%%%%%%%%%
%%%%%%%%%%%%%%%%%%%%%%%%%%%%%%%%%%%%%%%%%%%%%%%%
\section{Finite Fourier Transform and Convolution}
\label{sec-dft}

%%%%%%%%%%%%%%%%%%%%%%%%%%%%%%%%%%%%%%%%%%%
\subsection{Discrete Orthonormal Bases}

A discrete signal is a finite-dimensional vector $f \in \CC^N$, where $N$ is the number of samples and $f_n$ is the value at a specified location. For a 2-D image $f \in \CC^N \simeq \CC^{N_0 \times N_0}$, we have $N = N_0 \times N_0$, with $N_0$ pixels in each direction.

The discrete inner product is the analogue of the continuous $\Ldeux$ inner product:
\eq{
	\dotp{f}{g} = \sum_{n=0}^{N-1} f_n \bar g_n.
}
The corresponding squared Euclidean distance is
\eq{
	\norm{f-g}^2 = \sum_{n=0}^{N-1} |f_n-g_n|^2.
}

Exactly as in the continuous case, a discrete orthonormal basis $\{ \psi_k \}_{0 \leq k < N }$ of $\CC^N$ satisfies
\eql{\label{eq-orthobasis-disc}
	\dotp{\psi_k}{\psi_{k'}} = \de_{k-k'}.
}
A signal expands in this orthonormal basis as
\eq{
	f = \sum_{k=0}^{N-1} \dotp{f}{\psi_k} \psi_k.
}
It preserves energy
\eq{
	\norm{f}^2 = \sum_{n=0}^{N-1} |f_n|^2 = 
\sum_{k=0}^{N-1} |\dotp{f}{\psi_k}|^2
}

Computing the set of all inner products $\{\dotp{f}{\psi_k}\}_{0 \leq k < N}$ directly requires $O(N^2)$ operations.
This cost is prohibitive when $N$ reaches millions of samples. Structured bases permit faster decompositions: Fourier and wavelet transforms require $O(N\log(N))$ and $O(N)$ operations, respectively.


%%%%%%%%%%%%%%%%%%%%%%%%%%%%%%%%%%%%%%%%%%%%%%%%
\subsection{Discrete Fourier transform}

We denote $f = (f_n)_{n=0}^{N-1} \in \RR^N$, and regard its entries as indexed by $n \in \ZZ/N\ZZ$, which is a finite abelian group under addition. This corresponds to using periodic boundary conditions.

The discrete Fourier transform is defined as
\eql{\label{eq-dft}
	\foralls k=0,\ldots,N-1, \quad \hat f_k \eqdef \sum_{n=0}^{N-1} f_n e^{-\piN kn} = \dotp{f}{\phi_k}
	\qwhereq
	\phi_k \eqdef ( e^{\piN kn} )_{n=0}^{N-1} \in \CC^N
}
where the canonical inner product on $\CC^N$ is $\dotp{u}{v}=\sum_{n=0}^{N-1}u_n\bar v_n$ for $(u,v) \in (\CC^N)^2$.
%
This definition is motivated by sampling the Fourier basis $x \mapsto e^{\imath k x}$ on $\RR/2\pi\ZZ$ at equally spaced points $( \frac{2\pi}{N}n )_{n=0}^{N-1}$.
%
The next proposition establishes the orthogonality of the discrete Fourier atoms and the associated reconstruction formula.

\begin{prop}
	One has
	\eq{
		\dotp{\phi_k}{\phi_\ell} = \choice{
			N \qifq k=\ell, \\
			0 \quad\text{otherwise.}
		}
	}
	In particular, this implies
	\eql{\label{eq-dft-inv}
		\foralls n=0,\ldots,N-1, \quad
		f_n = \frac{1}{N}\sum_k \hat f_k e^{\piN kn}. 
	}
\end{prop}
\begin{proof}
	One has, if $k \neq \ell$
	\eq{
		\dotp{\phi_k}{\phi_\ell} = \sum_n e^{\piN (k-\ell) n} = 
		\frac{ 1 - e^{\piN (k-\ell) N} }{ 1 - e^{\piN (k-\ell)} } = 0
	}
	which is the sum of a geometric series (equivalently, the sum of equispaced points on a circle). 
	%
	The inversion formula is simply $f = \sum_k \dotp{f}{\phi_k} \frac{\phi_k}{\norm{\phi_k}^2}$.
\end{proof}


%%%%%%%%%%%%%%%%%%%%%%%%%%%%%%%%%%%%%%%%%%%%%%%%
\subsection{Fast Fourier transform}

Assuming $N=2N'$, one has 
\begin{align*}
	\hat f_{2k} &= \sum_{n=0}^{N'-1} (f_n + f_{n+N/2}) e^{ -\frac{2\imath\pi}{N'} k n } \\
	\hat f_{2k+1} &= \sum_{n=0}^{N'-1} e^{ -\frac{2\imath\pi}{N} n } (f_n - f_{n+N/2}) e^{ -\frac{2\imath\pi}{N'} k n }.
\end{align*}
For the second line, we used the computation
\eq{
	e^{-\frac{2\imath\pi}{N}(2k+1)(n+N/2)} =
	e^{-\frac{2\imath\pi}{N}(2kn+kN+n+N/2)} = - e^{ -\frac{2\imath\pi}{N} n } e^{ -\frac{2\imath\pi}{N'} k n }.
}
Let $\Ff_N(f)=\hat f$ denote the discrete Fourier transform on $\CC^N$. Introduce
the vectors $f_e = (f_n + f_{n+N/2})_n \in \CC^{N'}$ and $f_o = (f_n - f_{n+N/2})_n \in \CC^{N'}$
and the twiddle factors $\al_N=(e^{-2\imath\pi n/N})_{n=0}^{N'-1}\in\CC^{N'}$. The recursion becomes
\eq{
	\Ff_N(f) = \Ii_N( \Ff_{N/2}(f_e), \Ff_{N/2}(f_o \odot \al_N) )
}
where $\Ii_N$ is the ``interleaving'' operator, defined by $\Ii_N(a,b) \eqdef (a_1,b_1,a_2,b_2,\ldots,a_{N'},b_{N'})$.
%
Recursing yields the radix-two fast Fourier transform (FFT) when $N$ is a power of two. Other FFT factorizations handle general lengths. Zero-padding to a power of two is also useful, although it changes the transform length and the sampled frequencies.

This algorithm can also be interpreted as a factorization of the Fourier matrix into a product of $O(\log N)$ sparse matrices.

\begin{figure}
\centering
\includegraphics[width=.5\linewidth]{2-fourier/fft}
\caption{\label{fig-fft}
Diagram of one step of the FFT.
}
\end{figure}
 

Let $C(N)$ denote the number of elementary operations required to compute $\hat f$. The recursion gives
\eql{\label{eq-cplx-fft}
	C(N) = 2C(N/2)+NK
}
where $K N$ accounts for $N$ complex additions and $N/2$ multiplications. With the change of variable
\eq{
	\ell \eqdef \log_2(N)
	\qandq 
	T(\ell) \eqdef \frac{C(N)}{N}
}
i.e. $C(N)=2^\ell T(\ell)$, the relation~\eqref{eq-cplx-fft} becomes
\eq{
	2^\ell T(\ell) = 2 \times 2^{\ell-1} T(\ell-1) + 2^\ell K \qarrq
	T(\ell) = T(\ell-1) + K \qarrq
	T(\ell) = T(0) + K \ell
} 
and using the fact that $T(0)=C(1)/1=0$, one obtains
\eq{
	C(N) = K N \log_2(N).
}
Direct evaluation instead requires $O(N^2)$ operations to compute the $N$ coefficients~\eqref{eq-dft}, each involving a sum of size $N$.


%%%%%%%%%%%%%%%%%%%%%%%%%%%%%%%%%%%%%%%%%%%%%%%%
\subsection{Finite convolution}

For $(f,g) \in (\RR^{N})^2$, one defines $f\star g \in \RR^N$ as
\eql{\label{eq-finite-convol}
	\foralls n=0,\ldots,N-1, \quad
	(f \star g)_n \eqdef \sum_{k=0}^{N-1} f_k g_{n-k} = \sum_{k+\ell=n} f_k g_\ell 
}
where $+$ and $-$ are interpreted modulo $N$ (vectors are defined on the group $\ZZ/N\ZZ$, or equivalently, one uses periodic boundary conditions).
%
This defines a commutative algebra structure $(\RR^N,+,\star)$, with neutral element the ``Dirac'' $\de_0 \eqdef (1,0,\ldots,0)^\top \in \RR^N$. The following proposition shows that $\Ff : f \mapsto \hat f$ is an algebra isomorphism and an isometry (up to a scaling by $\sqrt{N}$ of the norm) from $(\CC^N,+,\star)$ onto $(\CC^N,+,\odot)$ with neutral element $\ones_N=(1,\ldots,1) \in \RR^N$.

\begin{prop}\label{prop-tfd-conv}
	One has $\Ff( f \star g ) = \hat f \odot \hat g$. 
\end{prop}
\begin{proof}
	We denote $T : g \mapsto f \star g$. One has 
	\eq{
		(T \phi_\ell)_n = (f \star \phi_\ell)_n = \sum_p f_p e^{ \piN \ell(n-p) } = e^{ \piN \ell n } \hat f_\ell = \hat f_\ell (\phi_\ell)_n.
	}
	Thus $(\phi_\ell)_\ell$ form a basis of $N$ eigenvectors of $T$, with eigenvalues $\hat f_\ell$. The operator $T$ is diagonal in this basis. Let $F=(  e^{ -\piN k n } )_{k,n}$ be the Fourier transform matrix. The inversion formula~\eqref{eq-dft-inv} gives $F^{-1} = \frac{1}{N}F^*$, where $F^*=\bar F^\top$ is the adjoint (conjugate transpose). The diagonalization of $T$ is therefore
	\eq{
		T = F^{-1} \diag(\hat f) F \qarrq
		\Ff(T g) = \diag(\hat f) F g \qarrq
		\Ff(f \star g) = \diag(\hat f) \hat g.
	} 
\end{proof}

This proposition gives an $O(N\log N)$ convolution algorithm:
\eq{	
	f \star g = \Ff^{-1}( \hat f \odot \hat g ).
}
This is more efficient than directly evaluating formula~\eqref{eq-finite-convol} for filters with large support. When $|\Supp(g)|=P$ is small, direct evaluation costs $O(PN)$ and may be faster. An example is $g=[1,1,0,\ldots,0,1]/3$, the moving average, where
\eq{
	(f \star g)_n = \frac{f_{n-1}+f_n+f_{n+1}}{3}
}
needs $3N$ operations.


\paragraph{Convolution as translation invariant operator.}

Define translation on $\ZZ/N\ZZ$ by $(T_\tau f)_n \eqdef f_{n-\tau}$, with $n-\tau$ computed modulo $N$. The next proposition is the finite counterpart of the results on $\RR$ and $\RR/2\pi\ZZ$. On $\ZZ/N\ZZ$, no convergence issues arise, and the impulse response can be defined directly.

\begin{prop}
	Let $H:\RR^N\to\RR^N$ be linear. Then $H \circ T_\tau = T_\tau \circ H$ for all $\tau$ if and only if there exists $h \in \RR^N$ (often called the impulse response) so that $Hf = f \star h$.
\end{prop}

\begin{proof}
	For the convolution-to-invariance implication, we first note that $T_\tau f=f\star\delta_\tau$ is a convolution against a translated Dirac $\delta_\tau=T_\tau(\delta)$ where $\delta=(1,0,\ldots,0)$. If $Hf=f\star h$, commutativity gives
	$$
		H \circ T_\tau(f) = h \star \delta_\tau \star f = \delta_\tau \star h \star f
		= T_\tau \circ H(f).
	$$	
	Conversely, we define $h \eqdef H(\delta)$. Then
	$$
		H(f) = H( \sum_i f_i (T_{i} \delta) ) = \sum_i f_i H(T_i \delta) = \sum_i f_i T_i H(\delta) = \sum_i f_i h_{\cdot-i} = f \star h.
	$$
	 
\end{proof}


\paragraph{Polynomial multiplication.}


The FFT can multiply large polynomials. It can also multiply large integers by interpreting their digit expansions in a fixed base as polynomials. Indeed,
\eq{
	(\sum_{i=0}^A a_i X^i)(\sum_{j=0}^B b_j X^j) =
	\sum_{k=0}^{A+B} ( \sum_{i+j=k} a_i b_j ) X^k
}
Setting aside the carry operation, which can be performed in linear time,
one can write $\sum_{i+j=k} a_i b_j = (\bar a \star \bar b)_k$ when one defines $\bar a,\bar b\in\RR^{A+B+1}$ by zero padding.


%%%%%%%%%%%%%%%%%%%%%%%%%%%%%%%%%%%%%%%%%%%%%%%%
%%%%%%%%%%%%%%%%%%%%%%%%%%%%%%%%%%%%%%%%%%%%%%%%
%%%%%%%%%%%%%%%%%%%%%%%%%%%%%%%%%%%%%%%%%%%%%%%%
\section{Discretization Issues}
\label{sec-fourier-discretization}

The FFT also approximates the continuous Fourier transform and its inverse. Reversing the roles of space and frequency leads to an efficient spectral interpolation method.


%%%%%%%%%%%%%%%%%%%%%%%%%%%%%%%%%%%%%%%%%%%%%%%%
\subsection{Fourier approximation via spatial zero padding.}

The discrete Fourier transform~\eqref{eq-dft} approximates samples of the continuous Fourier transform~\eqref{eq-fourier-transform}. For a sufficiently smooth $f$ supported on $[0,1]$, consider the vector $f^Q \eqdef (f(n/N))_{n=0}^{Q-1} \in \RR^Q$. Taking $Q\geq N$ pads the original samples with zeros, since continuity and the support condition give $f(n/N)=0$ for $n\geq N$.
For signed frequency indices $-\lfloor Q/2\rfloor\leq k\leq\lceil Q/2\rceil-1$ and $T=Q/N$,
\begin{align*}
 \frac1N\hat f_k^Q
 &=\frac1N\sum_{n=0}^{N-1}f(n/N)e^{-2\pi\imath nk/Q}\\
 &\approx\int_0^1f(x)e^{-2\pi\imath kx/T}\,\d x
 =\hat f(2\pi k/T).
\end{align*}
For a fixed physical frequency, the Riemann-sum error is $O(1/N)$ for $f\in C^1$. The bound need not be uniform over frequencies growing with $N$. Increasing $Q$ at fixed $N$ samples the same spectrum more densely; it does not increase the Nyquist frequency.
% FAUX Increasing the amount $Q$ of zero padding is a way to compute larger frequencies. Increasing the discretization precision $N$ is on contrary a way to increase the precision of the Fourier sampling (using smaller step size $2\pi/T$). 

\begin{figure}
\centering
\includegraphics[width=.75\linewidth]{2-fourier/padding-spacial}
\caption{\label{fig-padding-spacial}
Fourier transform approximation by zero-padding in the spatial domain.
}
\end{figure}


%%%%%%%%%%%%%%%%%%%%%%%%%%%%%%%%%%%%%%%%%%%%%%%%
\subsection{Fourier interpolation via spectral zero padding.}

One can reverse the roles of space and frequency in the previous construction. 
%
Given $N$ uniform discrete samples $f^N=(f^N_n)_{n=0}^{N-1}$, one can compute their discrete Fourier transform $\Ff(f^N) = \hat f^N$ (in $O(N \log(N))$ operations with the FFT),
\eq{
	\hat f^N_k \eqdef \sum_{n=0}^{N-1} f^N_n e^{-\frac{2\imath\pi}{N} n k}, 
}
and then zero-pad the spectrum to obtain a vector of length $Q\geq N$.
%
For simplicity, assume that $N=2N'+1$ is odd. Even lengths admit a similar, slightly more involved construction.
%
With frequency indices $-N'\leq k\leq N'$, define the padded spectrum by
\eq{
	\tilde f_k^Q=\hat f_k^N\quad(-N'\leq k\leq N'),\qquad\tilde f_k^Q=0\quad\text{otherwise},\qquad\tilde f^Q\in\CC^Q
} 

The signed frequency indices are interpreted modulo $Q$ when calling an FFT routine.
An inverse transform of size $Q$, multiplied by $Q/N$, then gives in $O(Q\log Q)$ operations
\begin{align*}
	\frac{Q}{N} \Ff^{-1}( \tilde f^Q )_\ell &= 
	\frac{Q}{N} \times \frac{1}{Q} \sum_{k=-N'}^{N'} \hat f_k^N e^{ \frac{2\imath\pi}{Q} \ell k }
	= \frac{1}{N} \sum_{k=-N'}^{N'} \sum_{n=0}^{N-1} f^N_n e^{-\frac{2\imath\pi}{N} n k} e^{ \frac{2\imath\pi}{Q} \ell k }  \\
	&= \sum_{n=0}^{N-1} f^N_n \frac{1}{N} \sum_{k=-N'}^{N'} e^{ 2\imath\pi \pa{ -\frac{n}{N} + \frac{\ell}{Q} } k } 
	=  \sum_{n=0}^{N-1} f^N_n 
			\frac{ 
				\sin\left[ \pi N \pa{ \frac{\ell}{Q} - \frac{n}{N} } \right]  
			}{
				N \sin\left[ \pi \pa{ \frac{\ell}{Q} - \frac{n}{N} } \right]
			}  \\
	&= \sum_{n=0}^{N-1} f^N_n \sinc_N\pa{ \frac{\ell}{T} - n }
	\qwhereq
	T \eqdef \frac{Q}{N} \qandq \sinc_N(u) \eqdef \frac{ \sin(\pi u) }{ N \sin(\pi u /N) }.
\end{align*}
We use the geometric-series identity with $\rho=e^{\imath \om}$, $a=-b$, $\om = 2\pi \pa{ -\frac{n}{N} + \frac{\ell}{Q} }$,
\eq{
	\sum_{i=a}^b \rho^i = \frac{ \rho^{a-\frac{1}{2}} - \rho^{b+\frac{1}{2}} }{ \rho^{-\frac{1}{2}} - \rho^{\frac{1}{2}} }
	= \frac{ \sin( (b+\frac{1}{2}) \om ) }{ \sin(\om/2) }.
}
This gives a discrete counterpart of the Shannon interpolation formula~\eqref{eq-shannong-interp}. It evaluates the trigonometric interpolant of the samples exactly on a grid of size $Q$, at cost $O(Q\log(Q))$. Furthermore, $\sinc_N(u)\to\sinc(u)$ for fixed $u$ as $N\to\infty$, with removable singularities filled by continuity.

\begin{figure}
\centering
\includegraphics[width=.75\linewidth]{2-fourier/padding-fourier}
\caption{\label{fig-padding-fourier}
Interpolation by zero-padding in the frequency domain.
}
\end{figure}

%%%%%%%%%%%%%%%%%%%%%%%%%%%%%%%%%%%%%%%%%%%%%%%%
%%%%%%%%%%%%%%%%%%%%%%%%%%%%%%%%%%%%%%%%%%%%%%%%
%%%%%%%%%%%%%%%%%%%%%%%%%%%%%%%%%%%%%%%%%%%%%%%%
\section{Fourier in Multiple Dimensions}
\label{sec-fourier-multiple-dim}

Tensor products extend the one-dimensional Fourier transform to any finite dimension $d>1$.


%%%%%%%%%%%%%%%%%%%%%%%%%%%%%%%%%%%%%%%%%%%%%%%%
\subsection{On Continuous Domains}

%%%
\paragraph{On \texorpdfstring{$\RR^d$}{Euclidean Space}.}

\wrapf{2-fourier/wave}{2-D sine wave.}
Tensor products retain a useful feature of Fourier atoms: a product of elementary 1-D complex exponentials is a plane wave
\eq{
	\prod_{\ell=1}^d e^{ \imath x_\ell \om_\ell } = e^{ \imath \dotp{x}{\om} }
}
whose constant-phase hyperplanes are orthogonal to the wave vector $\om=(\om_\ell)_{\ell=1}^d \in \RR^d$ when this vector is nonzero. Here $\dotp{x}{\om} = \sum_\ell x_\ell \om_\ell$ is the canonical inner product on $\RR^d$.

The Fourier transform and its inverse are defined by
\begin{align*}
	\foralls \om \in \RR^d, \quad \hat f(\om) &\eqdef \int_{\RR^d} f(x) e^{-\imath\dotp{x}{\om}} \d x, \\
	\foralls x \in \RR^d, \quad  f(x) &= \frac{1}{(2\pi)^d} \int_{\RR^d}\hat f(\om)e^{\imath\dotp{x}{\om}}\,\d\om,
\end{align*}
under the same integrability assumptions as in one dimension.


\myfigure{
\image{orthobases}{.3}{2d-extension-fourier-plane}
\image{orthobases}{.2}{2d-extension-fourier-1}
\image{orthobases}{.2}{2d-extension-fourier-2}
\image{orthobases}{.2}{2d-extension-fourier-3}
}{%
	2D Fourier orthogonal bases. %	
}{fig-2d-fourier-extension}


%%%
\paragraph{On \texorpdfstring{$(\RR/2\pi\ZZ)^d$}{the Torus}.}

Given an orthonormal basis $(\phi_{n_1})_{n_1 \in \NN}$ of $L^2(\XX)$, one constructs an orthonormal basis of $L^2(\XX^d)$ by tensorization
\eql{\label{eq-tensor-product}
	\foralls k=(k_1,\ldots,k_d) \in \NN^d, \quad
	\foralls x \in \XX^d, \quad
	\phi_k(x) = \phi_{k_1}(x_1) \ldots \phi_{k_d}(x_d).
}
Fubini's theorem gives orthogonality. Completeness gives convergence of the rectangular partial sums $\sum_{\norm{k}_\infty \leq N} \dotp{f}{\phi_k} \phi_k \rightarrow f$ in $L^2(\XX^d)$.

%%% FIG %%%
\wrapf{2-fourier/torus}{The two-dimensional torus $\TT^2=(\RR/2\pi\ZZ)^2$.}
%%%
For the multidimensional torus $(\RR/2\pi\ZZ)^d$, using the Fourier basis~\eqref{eq-fourier-basis-1d}, this gives the basis
\eq{
	\foralls k \in \ZZ^d, \quad \phi_k(x) = e^{\imath \dotp{x}{k}}
} 
which is indeed a Hilbertian orthonormal basis for the inner product $\dotp{f}{g} \eqdef \frac{1}{(2\pi)^d} \int_{\TT^d} f(x)\bar g(x)\d x$. 
%
This defines the Fourier transform and the reconstruction formula on $L^2(\TT^d)$
\eq{
	\hat f_k \eqdef \frac{1}{(2\pi)^d} \int_{\TT^d} f(x) e^{-\imath \dotp{x}{k}}  \d x
	\qandq
	f = \sum_{k \in \ZZ^d} \hat f_k e^{\imath \dotp{x}{k}}.
}

%%%%%%%%%%%%%%%%%%%%%%%%%%%%%%%%%%%%%%%%%%%%%%%%
\subsection{On Discrete Domains}
\label{sec-fft-2d}

%%%%
\paragraph{Discrete Fourier Transform.}

\wrapf{2-fourier/torus-discr}{Discrete 2-D torus.}
On a $d$-dimensional discrete domain of the form
\eq{
	n = (n_1,\ldots,n_d) \in \YY_d \eqdef \range{1,N_1} \times \ldots \times \range{1,N_d}
}
(we denote $\range{a,b} \eqdef \enscond{i \in \ZZ}{a \leq i \leq b}$)
of $N=N_1\ldots N_d$ points, with periodic boundary conditions, one defines an orthogonal basis $(\phi_k)_k$ by the same tensor product formula as~\eqref{eq-tensor-product} but using the 1-D discrete Fourier basis~\eqref{eq-dft}
\eql{\label{eq-exp-four-disc-multid}
	\foralls (k,n) \in \YY_d^2, \quad
	\phi_k(n) = \phi_{k_1}(n_1) \ldots \phi_{k_d}(n_d) 
	= \prod_{\ell=1}^d e^{\frac{2\imath\pi}{N_\ell} k_\ell n_\ell}
	= e^{ 2\imath\pi \dotp{k}{n}_{\YY_d} } 	
} 
where we used the (rescaled) inner product
\eql{\label{eq-inner-prod-multi-d}
	\dotp{k}{n}_{\YY_d} \eqdef \sum_{\ell=1}^d \frac{k_\ell n_\ell}{N_\ell}.
}
The pairing in~\eqref{eq-inner-prod-multi-d} acts on frequency and spatial index vectors. The basis $(\phi_k)_k$ is orthonormal for the signal inner product $N^{-1}\sum_n f_n\bar g_n$. As in one dimension, we define the Fourier coefficients without normalizing by $(N_\ell)_\ell$:
\begin{align*}
	\foralls k \in \YY_d, \quad 
	\hat f_k &\eqdef \sum_{n \in \YY_d} f_n e^{-2\pi\imath\dotp{k}{n}_{\YY_d}}, \\
	\foralls n \in \YY_d, \quad 
	f_n &= \frac{1}{N} \sum_{k \in \YY_d} \hat f_k e^{2\pi\imath\dotp{k}{n}_{\YY_d}}.
\end{align*}



%%%%
\paragraph{Fast Fourier Transform.}

We describe the algorithm for $d=2$; the same construction works in any finite dimension.
%
A fast algorithm for orthogonal decompositions in two 1-D bases $(\phi_{k_1}^{1})_{k_1=1}^{N_1}$, $(\phi_{k_2}^{2})_{k_2=1}^{N_2}$ also gives a fast decomposition in the tensor product basis
$(\phi_{k_1}^{1} \otimes \phi_{k_2}^{2})_{k_1,k_2}$: apply the algorithm first to the rows and then to the columns of the matrix $(f_n)_{n=(n_1,n_2)} \in \RR^{N_1 \times N_2}$. Either order gives the same result.
%
Indeed
\eq{
	\foralls k=(k_1,k_2), \quad
	\dotp{f}{\phi_{k_1}^{1} \otimes \phi_{k_2}^{2}}
	= \sum_{n=(n_1,n_2)} f_n\overline{\phi_{k_1}^{1}(n_1)}\,\overline{\phi_{k_2}^{2}(n_2)} =
	\sum_{n_1} \pa{
		\sum_{n_2} f_{n_1,n_2}\overline{\phi_{k_2}^{2}(n_2)}
	}\overline{\phi_{k_1}^{1}(n_1)}.
}
If $C(N_1)$ is the cost of the 1-D algorithm on $\RR^{N_1}$, the resulting 2-D transform costs $N_2 C(N_1)+N_1 C(N_2)$. For the FFT, this is $O(N_1 N_2 \log(N_1 N_2))=O(N \log(N))$, where $N=N_1 N_2$.

Represent $f \in \RR^{N_1 \times N_2}$ as a matrix, and write $F_N=(e^{-2\imath\pi kn/N})_{k,n}$ for the Fourier matrix, whose rows are the $\phi_k^*$. The 2-D transform can then be expressed as matrix products:
\eq{
	\hat f = F_{N_1} \times f \times F_{N_2}^{\top}\in\CC^{N_1\times N_2}.
}
In practice, the FFT performs these operations without explicit matrix multiplication.

\myfigure{
\tabquatre{
\image{fourier}{.238}{fourier-2d-image}&
\image{fourier}{.238}{fourier-2d-fft}&
\image{fourier}{.238}{fourier-2d-masked}&
\image{fourier}{.238}{fourier-2d-masked-fft}
}
}{%
	2-D Fourier analysis of an image (left), and attenuation of the periodicity artifact using masking (right). %
}{fig-fourier-2d}


\texttt{Associated code: coding/test\_fft2.m}


%%%%%%%%%%%%%%%%%%%%%%%%%%%%%%%%%%%%%%%%%%%%%%%%
\subsection{Shannon sampling theorem.}

Theorem~\ref{thm-shannon-sampling} extends to sampling on a uniform Cartesian grid in $\RR^d$ by tensorization. Use the continuous representative of the bandlimited function $f$. In two dimensions, if $\supp(\hat f) \subset [-\pi/s_1,\pi/s_1] \times [-\pi/s_2,\pi/s_2]$ and $f$ decays sufficiently fast,
\eq{
		\foralls x \in \RR^2, \quad 
		f(x) = \sum_{n \in \ZZ^2} f(n_1 s_1,n_2 s_2) \sinc(x_1/s_1-n_1) \sinc(x_2/s_2-n_2) \qwhereq
		\sinc(u) = \frac{\sin(\pi u)}{\pi u}.
}
	

%%%%%%%%%%%%%%%%%%%%%%%%%%%%%%%%%%%%%%%%%%%%%%%%
\subsection{Convolution in higher dimension.}

Convolution on $\XX^d$, for $\XX=\RR$ or $\XX=\RR/2\pi\ZZ$, is defined as in one dimension:
\eq{
	f \star g(x) = \int_{\XX^d} f(t) g(x-t) \d t. 
}
Similarly, finite discrete convolution of vectors $f \in \RR^{N_1 \times N_2}$ extends formula~\eqref{eq-finite-convol} as
\eq{
	\foralls n \in \range{0,N_1-1} \times \range{0,N_2-1}, \quad
	(f \star g)_n \eqdef \sum_{k_1=0}^{N_1-1} \sum_{k_2=0}^{N_2-1} f_k g_{n-k}  
}
where additions and subtractions of vectors are performed modulo $(N_1,N_2)$.

The Fourier-convolution identity $\Ff(f \star g)=\hat f \odot \hat g$ holds in each of these settings. In the finite case, it yields an $O(N \log(N))$ convolution algorithm even when $f$ and $g$ have large support.


%%%%%%%%%%%%%%%%%%%%%%%%%%%%%%%%%%%%%%%%%%%%%%%%
%%%%%%%%%%%%%%%%%%%%%%%%%%%%%%%%%%%%%%%%%%%%%%%%
%%%%%%%%%%%%%%%%%%%%%%%%%%%%%%%%%%%%%%%%%%%%%%%%
\section{Application to ODEs and PDEs}
\label{sec-fourier-pdes}

%%%%%%%%%%%%%%%%%%%%%%%%%%%%%%%%%%%%%%%%%%%%%%%%
\subsection{On Continuous Domains}

We give the main calculations, leaving analytic details of existence and regularity aside.

On $\XX=\RR$ or $\TT$, one has
\eq{
	\Ff( f^{(k)} )(\om) = (\imath \om)^k \hat f(\om). 
}
Intuitively, $f^{(k)} = f \star \de^{(k)}$ where $\de^{(k)}$ is a distribution with Fourier transform $\hat \de^{(k)}(\om) = (\imath \om)^k$. 
%
Similarly on $\XX=\RR^d$ (see Section~\ref{sec-fourier-multiple-dim} for the definition of the Fourier transform in dimension $d$), one has
\eql{\label{eq-lapl-fourier}
	\Ff(\Delta f)(\om) = -\norm{\om}^2 \hat f(\om)
}
The same identity holds on $\TT^d$, with $\om$ replaced by $n \in \ZZ^d$.
%
Fourier transforms and Fourier coefficients provide a powerful tool for studying linear differential equations with constant coefficients by turning them into algebraic equations.

\begin{figure}
\centering
\includegraphics[width=.5\linewidth]{2-fourier/heat}
\caption{\label{fig-heat}
Heat diffusion as a convolution.
}
\end{figure}

As a typical example, we consider the heat equation
\eq{
	\frac{\partial f_t}{\partial t} = \Delta f_t
	\qarrq
	\foralls \om, \quad \frac{\partial \hat f_t(\om)}{\partial t} = -\norm{\om}^2 \hat f(\om).
}
Hence $\hat f_t(\om) = \hat f_0(\om) e^{-\norm{\om}^2 t}$. Applying the inverse Fourier transform and the convolution theorem gives
\eq{
	f_t = G_t \star f_0 \qwhereq
	G_t=\frac{1}{(4\pi t)^{d/2}} e^{-\frac{\norm{x}^2}{4t}}
}
On $\RR^d$, $G_t$ is a Gaussian density with standard deviation $\sqrt{2t}$ in each direction. On the torus, periodize this kernel and adjust its mass to the chosen Haar normalization.


%%%%%%%%%%%%%%%%%%%%%%%%%%%%%%%%%%%%%%%%%%%%%%%%
\subsection{Finite Domain and Discretization}

On $\ZZ/N\ZZ$, corresponding to a discrete periodic domain, consider the forward first difference and the centered second difference:
\begin{align}\label{eq-disc-diff-1}
	D_1 f &\eqdef N(f_{n+1}-f_n)_n=f\star d_1\qwhereq d_1=N[-1,0,\ldots,0,1]^\top \in \RR^N, \\
	\label{eq-disc-diff-2}
	D_2 f=-D_1^\top D_1 f & \eqdef N^2 (f_{n+1}+f_{n-1}-2f_n)_n = f \star d_2 \qwhereq d_2=-d_1\star\bar d_1=N^2[-2,1,0,\ldots,0,1]^\top \in \RR^N.
\end{align}
\wrapf{2-fourier/finite-diff}{Comparison of the spectra of $\De$ and $D_2$.}
Thanks to Proposition~\ref{prop-tfd-conv}, one can equivalently compute
\eql{\label{eq-fft-lapl}
	\Ff( D_2 f ) = \hat d_2 \odot \hat f
	\qwhereq
	(\hat d_2)_k = N^2 ( e^{\piN k}+e^{-\piN k} - 2 ) = -4N^2 \sin\pa{\frac{\pi k}{N}}^2.
}
Here $\bar d_1[n]=d_1[-n]$ denotes periodic reflection. For fixed nonzero $k$ and $N\to\infty$, one has $(\hat d_2)_k \sim -(2\pi k)^2$ which matches the scaling of~\eqref{eq-lapl-fourier}.


%%%%%%%%%%%%%%%%%%%%%%%%%%%%%%%%%%%%%%%%%%%%%%%%
%%%%%%%%%%%%%%%%%%%%%%%%%%%%%%%%%%%%%%%%%%%%%%%%
%%%%%%%%%%%%%%%%%%%%%%%%%%%%%%%%%%%%%%%%%%%%%%%%
\section{A Bit of Group Theory}

The reference for this section is~\cite{peyre2004algebre}.

%%%%%%%%%%%%%%%%%%%%%%%%%%%%%%%%%%%%%%%%%%%%%%%%
\subsection{Characters}

For $(G,+)$ a commutative group, a character is a group morphism $\chi : (G,+) \rightarrow (\CC^*,\cdot)$, i.e. it satisfies
\eq{
	\foralls (n,m) \in G, \quad \chi(n+m) = \chi(n)\chi(m).
}
The characters form the dual group $(\hat G,\odot)$ under pointwise multiplication $(\chi_1 \odot \chi_2)(n) \eqdef \chi_1(n) \chi_2(n)$. The inverse of a character $\chi$ is $\chi^{-1}(n)=\chi(-n)$.

For a finite group $G$ with $|G|=N$, we have $N \times n=0$ for every $n \in G$. Consequently, $\chi(n)^N=\chi(N n)=\chi(0)=1$: characters take values among the roots of unity,
\eql{\label{eq-char-finite-group}
	\chi(n) \in \enscond{ e^{\piN k} }{ 0 \leq k \leq N-1 }.
}
Thus $\hat G$ is a finite group (since there are finitely many maps between two finite sets) and $\chi^{-1}=\bar\chi$. In the case of a cyclic group, the dual is simple to describe.

\begin{prop}\label{prop-iso-cyclic}
	For $G=\ZZ/N\ZZ$, the dual is $\hat G=( \phi_k )_{k=0}^{N-1}$, where $\phi_k=( e^{\piN nk} )_n$. The map $k \mapsto \phi_k$ defines a noncanonical isomorphism $G \sim \hat G$.
\end{prop}
\begin{proof}
	The $\phi_k$ are indeed characters.
	
	Conversely, for any $\chi \in \hat G$, according to~\eqref{eq-char-finite-group}, $\chi(1)=e^{\piN k}$ for some $k$. Then 
	\eq{
		\chi(n)=\chi(1)^n=e^{\piN k n} = \phi_k(n).
	}
	These characters are distinct because the values $\phi_k(1)$ are distinct. Hence $|G|=|\hat G|$, and the displayed map is an isomorphism.
\end{proof}

This proposition thus shows that characters of cyclic groups are exactly the orthogonal discrete Fourier basis defined in~\eqref{eq-dft}.


%%%
\paragraph{Commutative groups.}

The structure theorem for finitely generated abelian groups expresses every such group as a product of cyclic groups:
\eql{\label{eq-structure}
	G \sim (\ZZ/N_1\ZZ) \times \ldots \times (\ZZ/N_d\ZZ) \times \ZZ^Q.
}
If $G$ is finite, then $Q=0$ and $N=N_1\cdots N_d$. In this case, $G$ is simply a discrete $d$-dimensional ``rectangle'' with periodic boundary conditions.

For two finite groups $(G_1,G_2)$ one has 
\eql{\label{eq-iso-product-char}
	\widehat{G_1 \times G_2} = \hat G_1 \otimes \hat G_2 = \enscond{ \chi_1 \otimes \chi_2 }{ (\chi_1,\chi_2) \in \hat G_1 \times \hat G_2 }.
} 
Here $\otimes$ is the tensor product of two functions
\eq{
	\foralls (n_1,n_2) \in G_1 \times G_2, \quad
	(\chi_1 \otimes \chi_2)(n_1,n_2) \eqdef \chi_1(n_1)\chi_2(n_2).
}
The tensor product $\chi_1 \otimes \chi_2$ is a character. Conversely, every character has the factorization $\chi=\chi(\cdot,0) \otimes \chi(0,\cdot)$, since $(n_1,n_2)=(n_1,0)+(0,n_2)$.

Combining this factorization with the structure theorem proves the following isomorphism.

\begin{prop}
	If $G$ is commutative and finite then $\hat G \sim G$.
\end{prop}
\begin{proof}
	The structure theorem~\eqref{eq-structure} with $Q=0$ and the product-duality identity~\eqref{eq-iso-product-char} give
	\eq{
		\hat G \sim \hat G_1 \otimes \ldots \otimes \hat G_d
	} 
	where $G_\ell \eqdef \ZZ/N_\ell\ZZ$.
	%
	One then observes that  $\hat G_1 \otimes \hat G_2 \sim \hat G_1 \times \hat G_2$.
	%
	Proposition~\ref{prop-iso-cyclic} then completes the proof, since $\hat G_k \sim G_k$.
\end{proof}

The isomorphism $\hat G \sim G$ is noncanonical: it depends on a choice of generators and corresponding roots of unity. As in vector-space duality, the double-dual isomorphism $\hat{\hat{G}} \sim G$ is canonical and is given by evaluation:
\eq{
	g \in G \longmapsto e_g \in \hat{\hat{G}}
	\qwhereq
	\pa{
	e_g : \chi \in \hat G \mapsto \chi(g) \in \CC^*.
	}
}

%%%
\paragraph{Discrete Fourier transform from the Character Viewpoint.}

This construction also identifies the Fourier basis explicitly. The characters in $\hat G_\ell$ are the discrete Fourier atoms~\eqref{eq-dft}, of the form
\eq{
	( e^{\frac{2\imath\pi}{N_\ell} k_\ell n_\ell})_{n_\ell=0}^{N_\ell-1}
	\quad\text{for some}\quad
	0 \leq k_\ell < N_\ell.
} 
Identifying $G$ and $G_1 \times \ldots \times G_d$, 
their tensor products show that the characters in $\hat G$ form exactly the orthogonal multidimensional Fourier basis~\eqref{eq-exp-four-disc-multid}.

%%%%%%%%%%%%%%%%%%%%%%%%%%%%%%%%%%%%%%%%%%%%%%%%
\subsection{More General cases}

%%%
\paragraph{Infinite groups.}

For an infinite finitely generated abelian group, the structure theorem has $Q>0$. For locally compact abelian groups, the dual consists of continuous characters with values on the unit circle; invariant Haar measure determines integration. For $G=\ZZ$, the characters have a continuous frequency parameter,
\eq{
	\hat \ZZ = \enscond{ \phi_\om : n \mapsto e^{\imath n \om} \in \CC^\ZZ }{  \om \in \RR/2\pi\ZZ }
}
so that $\hat \ZZ \sim \RR/2\pi\ZZ$.
%
The case $G=\ZZ^Q$ follows by tensorization. 
%
The $(\phi_\om)_\om$ are ``orthogonal'' in the sense that $\dotp{\phi_\omega}{\phi_{\omega'}}_\ZZ=2\pi\delta_{2\pi}(\omega-\omega')$ can be understood as a Dirac kernel (this is similar to the Poisson formula), where $\dotp{u}{v}_\ZZ \eqdef \sum_n u_n \bar v_n$.
%
Expanding a sequence $(c_n)_{n \in \ZZ}$ in characters amounts to forming the Fourier series $\sum_n c_n e^{-\imath n \om}$.

Similarly, for $G=\RR/2\pi\ZZ$, one has $\hat G=\ZZ$, with orthonormal characters $\phi_n=e^{\imath \cdot n}$, so that the decomposition of functions in $L^2(G)$ is the computation of Fourier coefficients.
%
On a compact group, normalized invariant measure provides the natural $L^2$ inner product.

%%%
\paragraph{Non-commutative groups.}
 
For a noncommutative group, one-dimensional characters do not recover $G$. For example, the symmetric group $\Si_N$ on $N\geq2$ letters has only the characters $\hat G=\{1,\epsilon\}$, where $\epsilon(\si) = (-1)^{q}$ is the signature and $q$ is the number of transpositions in a decomposition of $\si \in \Si_N$.

To study noncommutative groups, replace homomorphisms $\chi : G \rightarrow \CC^*$ by homomorphisms $\rho : G \rightarrow \text{GL}(\CC^{n_\rho})$ for some $n_\rho$. These are called representations of $G$. For $(g,g') \in G$, writing the group operation multiplicatively on $G$, the defining identity is $\rho(gg')=\rho(g) \circ \rho(g')$. When $n_\rho=1$, the identification $\text{GL}(\CC) \sim \CC^*$ recovers the one-dimensional characters. For any representation $\rho$, the function $\chi(g) \eqdef \tr(\rho(g))$, where $\tr$ denotes trace, is called its character. This trace function is generally not a group homomorphism.

For a finite group $G$, every subspace $V$ invariant under all $\rho(g)$ admits an invariant complement $W$, so that $\CC^{n_\rho}=V\oplus W$ with $W$ also invariant. In a basis adapted to these subspaces, the representation matrices are block diagonal. To construct the complement, use the inner product
\eq{
	\dotp{x}{y} \eqdef \sum_{g \in G} \dotp{\rho(g)x}{\rho(g)y}_{\CC^{n_\rho}}
}
and take the orthogonal complement $V^\bot$ for this inner product. In an orthonormal basis of $\CC^{n_\rho}$ for this product, all matrices $\rho(g)$ are unitary.
%
A representation is irreducible if it has no nonzero proper invariant subspace. These are the elementary representations from which larger ones are built by block-diagonal assembly.
%
We identify irreducible representations up to isomorphism: $(\rho,\rho')$ are isomorphic if $\rho'(g)=U^{-1}\rho(g)U$ for every $g$, with a single change-of-basis matrix $U \in \text{GL}(\CC^d)$. Their isomorphism classes form the unitary dual, denoted $\hat G$, which is generally a set rather than a group.

For the symmetric group, there is an explicit description of the set of irreducible representations.  
%
For example, $G=\Sigma_3$ has $|G|=6$ elements. It has two one-dimensional representations, the trivial representation and the signature, and one two-dimensional irreducible representation obtained by identifying $G$ with the isometries of an equilateral triangle in $\RR^2$. Their dimensions satisfy $6 = |G| = \sum n_\rho^2 = 1+1+2^2$.

The dimensions $n_\rho$ of the irreducible representations $\rho\in\hat G$ satisfy $\sum_{\rho \in \hat G} n_\rho^2 = N$. After choosing unitary representatives, their matrix entries form an orthogonal basis for functions $f : G \rightarrow \CC$. This basis is noncanonical because it depends on a choice of basis within each representation.
%
The associated Fourier transform of a function $f:G\rightarrow\CC$ is defined as
\eq{
	\foralls \rho \in \hat G, \quad
	\hat f(\rho) \eqdef \sum_{g \in G} f(g) \rho(g) \in \CC^{n_\rho \times n_\rho}.  
}
With this sign convention, these are pairings with conjugated matrix entries. Their orthogonality makes the transform invertible.
%
The next proposition gives the inverse transform.

\begin{prop}
One has
\eq{
	\foralls g \in G, \quad
	f(g) = \frac{1}{|G|} \sum_{\rho \in \hat G} n_\rho \tr( \hat f(\rho) \rho(g^{-1}) ). 
}
\end{prop}
\begin{proof}
The proof uses the regular-character identity, valid for every $g\in G$
\eq{
	A_g \eqdef \sum_{\rho \in \hat G} n_\rho \tr(\rho(g))=|G|\de_g
}
where $\de_g=0$ for $g \neq \Id_G$ and $\de_{\Id_G}=1$. A proof is given in the book by Diaconis, p.~13: decompose the character of the regular representation into irreducible characters and use their orthogonality. The regular representation acts by permuting, according to $G$, the canonical basis of $\RR^{|G|}$.
%
One then has 
\begin{align*}
	\frac{1}{|G|} \sum_{\rho \in \hat G} n_\rho \tr( \hat f(\rho) \rho(g^{-1}) ) 
	&=  \frac{1}{|G|} \sum_{\rho \in \hat G} n_\rho \tr( \sum_u f(u) \rho(u) \rho(g^{-1}) ) \\
	&=  \sum_u  f(u) \frac{1}{|G|}  \sum_{\rho \in \hat G} n_\rho \tr( \rho(u g^{-1}) )
	=  \sum_u f(u)A_{ug^{-1}}/|G| = f(g).
\end{align*}
\end{proof}

One can define the convolution of two functions $f,h : G \rightarrow \CC$ as
\eq{
	(f \star h)(a) \eqdef \sum_{bc=a} f(b) h(c)
}
Convolution need not be commutative on this group. The Fourier transform block-diagonalizes the convolution operators, as stated next.

\begin{prop}
Denoting $\Ff(f) = \hat f$
\eq{
	\Ff(f \star h)(\rho) = \hat f(\rho) \times \hat h(\rho)
}
where $\times$ denotes matrix multiplication.
\end{prop}
\begin{proof}
\eq{
	\Ff(f \star h)(\rho) = \sum_x \sum_y f(y) h(y^{-1}x) \rho(y y^{-1} x) = \sum_x \sum_y f(y) \rho(y) h(y^{-1}x) \rho(y^{-1}x) 
	= \sum_y f(y) \rho(y) \sum_z h(z) \rho(z)
}
where we made the change of variable $z=y^{-1}x$.
\end{proof}


Fast algorithms for computing $\hat f$ exist for some groups, including permutation groups. Their structure is considerably more involved than for $G=\ZZ/N\ZZ$.

For infinite compact groups, the theory uses a discrete infinite family of irreducible representations. The rotation group SO$(3)$ is a useful example. Its representations can be described explicitly through the spherical harmonics introduced in Section~\ref{sec-spherical-harm}. There is one representation of dimension $2\ell+1$ for each frequency index $\ell$.
 
%
% It is possible to lift the 3-D sphere $\SS^2$ on the group SO$(3)$ since it operate transitively on the sphere, for instance defining for $f : \SS^2 \rightarrow \CC$ its lifted version $f_L(g) = f(g(o))$ where $o$ is any point on the sphere. 

%%%%%%%%%%%%%%%%%%%%%%%%%%%%%%%%%%%%%%%%%%%%%%%%
%%%%%%%%%%%%%%%%%%%%%%%%%%%%%%%%%%%%%%%%%%%%%%%%
%%%%%%%%%%%%%%%%%%%%%%%%%%%%%%%%%%%%%%%%%%%%%%%%
\section{A Bit of Spectral Theory}

On a general domain $\XX$, the group-theoretic approach applies when $\XX=G$ is a group or when a group acts transitively on $\XX$. Another approach defines Fourier-like basis functions as eigenfunctions of a differential operator, following the characterization in Section~\ref{sec-fourier-pdes}. The Laplacian is particularly useful: it is second order, rotation invariant, and has natural counterparts on surfaces and graphs.


%%%%%%%%%%%%%%%%%%%%%%%%%%%%%%%%%%%%%%%%%%%%%%%%
\subsection{On a Surface or a Manifold}

\wrapf{2-fourier/laplacian-surf}{Computing the Laplacian on a surface.}
Let $\XX$ be a smooth compact connected Riemannian manifold of dimension $d$, without boundary. The Laplace--Beltrami operator can be characterized, for smooth $f$, by the small-ball mean-value expansion
\eq{
 (\Delta f)(x)=\lim_{\epsilon\to0}\frac{2(d+2)}{\epsilon^2}
 \left(\frac{1}{\mathrm{Vol}(B_\epsilon(x))}\int_{B_\epsilon(x)}f(y)\,\d\mu(y)-f(x)\right).
}
Here $\mu$ is Riemannian volume and $B_\epsilon(x)$ is the geodesic ball centered at $x$. The factor $\epsilon^{-2}$ is essential.

The operator $\Delta$ is unbounded and self-adjoint on its natural domain in $L^2(\XX)$; its resolvent is compact. It has a complete orthonormal family of eigenfunctions $(\phi_n)_{n\geq0}$, with eigenvalues, counted with multiplicity,
\eq{0=\lambda_0>\lambda_1\geq\lambda_2\geq\cdots\longrightarrow-\infty.}
The constant eigenfunction is $\phi_0=\mathrm{Vol}(\XX)^{-1/2}$, and the inner product is $\dotp{f}{g}_\XX=\int_\XX f(x)\overline{g(x)}\,\d\mu(x)$. Boundary conditions must be specified if the manifold has a boundary.

On the flat unit torus $\XX=(\RR/\ZZ)^d$,
\eq{
 \Delta f=\sum_{s=1}^d\frac{\partial^2f}{\partial x_s^2},
 \qquad\phi_n(x)=e^{2\pi\imath\dotp{n}{x}},
 \qquad\lambda_n=-4\pi^2\norm{n}^2,\quad n\in\ZZ^d.
}

%%%%%%%%%%%%%%%%%%%%%%%%%%%%%%%%%%%%%%%%%%%%%%%%
\subsection{Spherical Harmonics}
\label{sec-spherical-harm}

\wrapf{2-fourier/spherical}{Spherical coordinates.}
Consider the $(d-1)$-dimensional sphere $\SS^{d-1} = \enscond{x \in \RR^d}{\norm{x}_{\RR^d}=1}$.
% 
The Laplacian has explicit eigenfunctions. For $d=3$, they are indexed by $n=(\ell,m)$:
\eq{
	\foralls \ell \in \NN, \quad
	\foralls m=-\ell,\ldots,\ell, \quad
	\phi_{\ell,m}(\th,\phi) = e^{\imath m \phi} P_{\ell}^m( \cos(\th) )
}
The corresponding eigenvalue is $\la_{\ell,m} = -\ell(\ell+1)$. Here $P_{\ell}^m$ are associated Legendre functions (polynomials multiplied by $(1-x^2)^{|m|/2}$). Normalization makes the eigenfunctions orthonormal. We use spherical coordinates $x=(\sin\th\cos\phi,\sin\th\sin\phi,\cos\th)\in\mathbb S^2$ for $(\th,\phi) \in [0,\pi] \times [0,2\pi]$.
%
The index $\ell$ plays the role of the magnitude of a two-dimensional Fourier frequency.
%
For a fixed $\ell$, the space $V_\ell = \text{span}( \phi_{\ell,m} )$ is an eigenspace of $\Delta$, and is also invariant under rotation. 
	

%%%%%%%%%%%%%%%%%%%%%%%%%%%%%%%%%%%%%%%%%%%%%%%%
\subsection{On a Graph}

\wrapf{2-fourier/graph}{Weighted graph.}
Let $\XX$ be a graph with $N$ vertices indexed by $\{1,\ldots,N\}$. Its geometry is encoded by the weight matrix $W=(w_{i,j})_{1\leq i,j\leq N}$. The notation $i \sim j$ means that $(i,j)$ is an edge, for $(i,j) \in \XX^2$. Set the weights to zero on nonedges and assume nonnegative, symmetric weights: $w_{i,j}=w_{j,i}$.

The graph Laplacian $\Delta:\RR^N\to\RR^N$ sums weighted differences between neighboring values
\eq{
	\foralls f \in \RR^N, \quad
	(\Delta f)_i \eqdef \sum_{j \sim i} w_{i,j} f_j - (\sum_{j \sim i} w_{i,j}) f_i 
	\qarrq 
	\Delta = W-D
}
where $D \eqdef \diag_i(\sum_{j \sim i} w_{i,j})$. In particular, note that $\Delta\ones=0$.

For instance, if $\XX = \ZZ/N\ZZ$ with the graph $i \sim i-1$ and $i \sim i+1$ (modulo $N$), then $\De$ is the finite difference Laplacian operator $\Delta=N^{-2}D_2$ for unit edge weights defined in~\eqref{eq-disc-diff-2}. This extends to any dimension by tensorization.

\begin{prop}
Define the graph-gradient operator $G : f \in \RR^N \mapsto ( \sqrt{w_{i,j}} (f_i-f_j) )_{(i,j)\in E}$, with one component per undirected edge. Then
\eq{
	-\Delta = G^\top G
	\qarrq
	\foralls f \in \RR^N, \quad 
	\dotp{ \Delta f }{f}_{\RR^N} = -\dotp{Gf}{Gf}_{\RR^{P}}.
}
where $P$ is the number of undirected edges $E=\enscond{(i,j)}{i \sim j, i<j}$.
\end{prop}
\begin{proof}
By symmetry of $W$ and counting each undirected edge once,
\begin{align*}
 \norm{Gf}^2
 &=\sum_{i<j}w_{ij}(f_i-f_j)^2\\
 &=\sum_i f_i^2\sum_jw_{ij}-\sum_{i,j}w_{ij}f_if_j\\
 &=\dotp{Df}{f}-\dotp{Wf}{f}=-\dotp{\Delta f}{f}.
\end{align*}
Equality of these quadratic forms proves $G^\top G=-\Delta$.
\end{proof}

Thus $\Delta$ is symmetric negative semidefinite and has an orthonormal eigenbasis. One can take $\phi_1=N^{-1/2}\ones$, with $0=\lambda_1\geq\lambda_2\geq\cdots\geq\lambda_N$. The multiplicity of zero is the number of connected components of the graph formed by the positive-weight edges; for a connected graph with $N>1$, $\lambda_2<0$. On a uniform periodic grid, the Fourier basis is an eigenbasis.

The quadratic energy $\sum_{i<j}w_{ij}(f_i-f_j)^2$ is the discrete analogue of the Dirichlet energy and provides a natural measure of smoothness on a graph.



%%%%%%%%%%%%%%%%%%%%%%%%%%%%%%%%%%%%%%%%%%%%%%%%
\subsection{Further Applications}

Polynomial multiplication provides another application (\texttt{Code text\_polymult.m}).
%
Analogous transforms can also be defined over suitable finite fields.
